\subsection{\texorpdfstring{\(p\)-环}{p-环}}

定义 1．——若环 C 的理想 pC 是极大理想，且 C 关于 pC-进拓扑分离且完备，则称 C 为 p-环。

设 C 为一个环；若 \(p1_{C}\) 是幂零的且 C 的理想 pC 是极大理想，则 C 是 p-环，因为 C 的 pC-进拓扑是离散的。特别地，每个特征为 p 的域都是 p-环。

命题 1．——设 C 为一个 p-环。

\begin{enumerate}
\def\labelenumi{\alph{enumi})}
\item
  环 C 是局部环，极大理想为 pC。
\item
  假设 \(p1_{\mathbb{C}}\) 是幂零的。设 \(d\) 是满足 \(p^{d}1_{\mathbb{C}} = 0\) 的最小正整数。\(\mathbb{C}\) 的理想形如 \(p^{k}\mathbb{C}\)，其中 \(0 \leqslant k \leqslant d\)，且有 \(p^{k}\mathbb{C} \neq p^{l}\mathbb{C}\)，这里 \(k\)、\(l\) 是满足 \(0 \leqslant k \leqslant d, 0 \leqslant l \leqslant d\) 的两个不同整数。\(\mathbb{C}\)-模 \(\mathbb{C}\) 的长度为 \(d\)。
\item
  假设 \(p1_{\mathbb{C}}\) 不是幂零的。则 C 是离散赋值环，其剩余域的特征为 \(p\)，其分式域的特征为 0。形如 \(p^{n}\mathbf{C}\)（\(n \in \mathbf{N}\)）的理想两两不同；它们构成 C 的全部非零理想。C-模 C 的长度不是有限的。
\end{enumerate}

断言 \(a)\) 由 III，第 2 节第 13 号命题 19 得到。

由假设，\(\bigcap_{n\geqslant 0}p^{n}\mathbf{C} = \{0\}\)。设 \(x\neq 0\) 属于 \(\mathbf{C}\)；存在整数 \(n\geqslant 0\)，使得 \(x\in p^n\mathbf{C}, x\notin p^{n + 1}\mathbf{C}\)；因而存在元素 \(y\)，它属于 \(\mathbf{C}\)，且 \(x = p^n y\)；由于 \(y\) 不属于 \(p\mathbf{C}\)，故 \(y\) 可逆。

假设 \(p1_{\mathbb{C}}\) 不是幂零的。若 \(x\)、\(x'\) 是 C 的两个非零元素，则存在两个整数 \(n \geqslant 0\)、\(n' \geqslant 0\) 以及 C 的两个可逆元素 \(y\)、\(y'\)，使得 \(x = p^n y\)、\(x' = p^{n'} y'\)。于是 \(xx' = p^{n+n'} yy' \neq 0\)，故 C 是整环。由于 C 是局部环但不是域，且 C 的极大理想 \(m_{\mathbb{C}} = p\mathbb{C}\) 是主理想，C 是离散赋值环（VI，第 3 节第 6 号命题 9）。于是根据上引文献命题 8，C 的非零理想形如 \(p^n\mathbb{C}\)，且两两不同。特别地，环 C 不是阿廷环，因而 C-模 C 的长度不是有限的。C 的剩余域 C/pC 的特征为 \(p\)。设 C 的分式域的特征为 \(q\)。有 \(p1_{\mathbb{C}} \neq 0\)，从而 \(p \neq q\)。此外，若 \(q\) 非零，则有 \(q1_{\mathbb{C}} = 0\)，从而 C/pC 的特征为 \(q \neq p\)，这是荒谬的。这就证明了 \(c\)。

假设 \(p1_{C}\) 是幂零的。设 \(d\) 是满足 \(p^{d}1_{C} = 0\) 的最小正整数。我们有理想序列

\[
\mathrm{C} \supset p \mathrm{C} \supset p ^ {2} \mathrm{C} \supset \dots \supset p ^ {d - 1} \mathrm{C} \supset p ^ {d} \mathrm{C} = \{0 \}.\tag{E}
\]

若 \(k\) 是满足 \(0 \leqslant k < d\) 且 \(p^k\mathbf{C} = p^{k+1}\mathbf{C}\) 的整数，则可推出

\[
p ^ {d - k - 1} p ^ {k} \mathrm{C} = p ^ {d - k - 1} p ^ {k + 1} \mathrm{C} = \{0 \}
\]

这与假设 \(p^{d-1}1_{\mathbf{C}} \neq 0\) 矛盾。因此，序列 (E) 的元素两两不同。设 \(\mathfrak{a}\) 是 C 的一个理想，令 \(k\) 为满足 \(\mathfrak{a} \supset p^k\mathbf{C}\) 的最小正整数。设 \(x\) 是非零元素，属于 \(\mathfrak{a}\)；我们已经看到，\(x\) 形如 \(p^m u\)，其中 \(m \geqslant 0\)，且 \(u\) 在 C 中可逆。因此 \(p^m\mathbf{C} \subset \mathfrak{a}\)，从而 \(m \geqslant k\)，最终 \(x \in p^k\mathbf{C}\)。总之，\(\mathfrak{a} = p^k\mathbf{C}\)。于是序列 (E) 是 C-模 C 的 Jordan–Hölder 序列，其长度为 \(d\)。

推论 1．——若 p-环 C 是整环，则它是离散赋值环，或是特征为 p 的域。

设 C 是整环。若 \(p1_{\mathbf{C}}\) 是幂零的，则 \(p1_{\mathbf{C}} = 0\)，而 \(\{0\}\) 是 C 的极大理想，故 C 是特征为 \(p\) 的域。若 \(p1_{\mathbf{C}}\) 不是幂零的，则由命题 1 的 c)，C 是离散赋值环。

推论 2．——设 C 为 p-环，且 a 是 C 中异于 C 的理想。环 C/a 是 p-环。

不妨设 \(\mathfrak{a} \neq \{0\}\)。于是存在整数 \(i \geqslant 1\)，使得 \(\mathfrak{a} = p^i\mathrm{C}\)；理想 \(p\mathrm{C} / \mathfrak{a}\) 是 \(\mathrm{C} / \mathfrak{a}\) 的极大理想，且 \(p^i1_{\mathrm{C} / \mathfrak{a}} = 0\)，故 \(\mathrm{C} / \mathfrak{a}\) 是一个 \(p\) -环。

设 C 为一个 \(p\) -环。称 C 的长度为 \(l(\mathbf{C})\)，并将其定义为 \(\overline{\mathbf{R}}\) 中所有满足 \(n \geqslant 1\) 且 \(p^{n-1}1_{\mathbf{C}} \neq 0\) 的整数所成集合的上确界。当 \(l(\mathbf{C})\) 有限时，它就是 C-模 C 的长度；当 \(l(\mathbf{C})\) 等于 \(+\infty\) 时，C-模 C 的长度不是有限的（命题 1）。

例．——1) 对每个整数 \(n \geqslant 1\)，环 \(\mathbf{Z}/p^{n}\mathbf{Z}\) 都是一个 \(p\)-环，长度为 \(n\)。环 \(\mathbf{Z}_{p}\) 是 \(p\)-进整数环，也是无限长度的 \(p\)-环。

\begin{enumerate}
\def\labelenumi{\arabic{enumi})}
\setcounter{enumi}{1}
\tightlist
\item
  设 K 为特征为 p 的完美域。根据第 1 节第 8 号命题 8，Witt 向量环 W(K) 是无限长度的 p-环。映射 \((a_{n})_{n\in\mathbf{N}} \mapsto a_{0}\) 通过取商诱导出从 W(K)/pW(K) 到域 K 的同构（上引文献命题 7）。对每个整数 \(n \geqslant 1\)，环
\end{enumerate}

\[
\mathrm{W} _ {n} (\mathrm{K}) = \mathrm{W} (\mathrm{K}) / p ^ {n} \mathrm{W} (\mathrm{K})
\]

是长度为 n 的 p-环。

命题 2．——设 C 和 C' 是两个 p-环，且 \(u\) 是从 C 到 C' 的同态。设 \(v\) 是从 \(\kappa_{\mathbf{C}} = \mathbf{C} / p\mathbf{C}\) 到 \(\kappa_{\mathbf{C}'} = \mathbf{C}' / p\mathbf{C}'\) 的同态，它由 u 通过取商诱导。

\begin{enumerate}
\def\labelenumi{\alph{enumi})}
\item
  有 \(l(\mathbf{C}) \geqslant l(\mathbf{C}')\)，且 u 为单射当且仅当 \(l(\mathbf{C}) = l(\mathbf{C}')\)。
\item
  u 为满射的充要条件是 v 为同构。
\item
  u 为同构的充要条件是 v 为同构且 \(l(\mathbf{C}) = l(\mathbf{C}')\)。
\end{enumerate}

令 \(n \geqslant 1\) 为整数。有 \(u(p^{n-1}1_{\mathbf{C}}) = p^{n-1}1_{\mathbf{C}'}\)，因此关系 \(p^{n-1}1_{\mathbf{C}'} \neq 0\) 蕴含 \(p^{n-1}1_{\mathbf{C}} \neq 0\)，并且当 \(u\) 为单射时与之等价。所以 \(l(\mathbf{C}') \leqslant l(\mathbf{C})\)，且当 \(u\) 为单射时等号成立。若 \(u\) 不是单射，则存在整数 \(i < l(\mathbf{C})\)，使得 \(u\) 的核是理想 \(p^i\mathbf{C}\)（属于 \(\mathbf{C}\)）；于是 \(p^i1_{\mathbf{C}'} = 0\)，从而 \(l(\mathbf{C}') \leqslant i\)。这就证明了 \(a\)。

由于 \(\kappa_{\mathrm{C}}\) 和 \(\kappa_{\mathrm{C}'}\) 是域，同态 \(v\) 是单射。若 \(u\) 是满射，则 \(v\) 也是满射，因而是同构。反之，设 \(v\) 是满射。于是对于每个整数 \(n \geqslant 0\)，映射 \(v_{n}: p^{n}\mathrm{C}/p^{n+1}\mathrm{C} \to p^{n}\mathrm{C}'/p^{n+1}\mathrm{C}'\) 由 \(u\) 得到且是满射。由于 C 关于 \(p\mathrm{C}\)-进滤子完备，而 \(\mathbf{C}'\) 关于 \(p\mathbf{C}'\)-进滤子分离，根据 III，§ 2，第 8 号的定理 1 的推论 2，\(u\) 是满射。这就证明了 \(b\)。

最后，c) 由 a) 和 b) 得出。

命题 3．——设 \((\mathbf{C}_{n}, \pi_{n,m})\) 是相对于指标集 N 的环的射影系统。假设 \(\mathbf{C}_{n}\) 对所有 \(n \in \mathbf{N}\) 都是 p-环，且同态 \(\pi_{n,m}\) 都是满射。则 \(C = \varprojlim C_{n}\) 是 p-环，并且对于所有 \(n \in \mathbf{N}\)，典范同态 \(\pi_{n}: C \to C_{n}\) 是满射，并诱导出从 \(\kappa_{C}\) 到 \(\kappa_{C_{n}}\) 的同构。

由于映射 \(\pi_{n,m}\) 是满射，映射 \(\pi_n\) 也具有同样性质（E，III，p. 58，命题 5）。我们来证明 C 是一个 \(p\)-环。令 \(d_n\) 为 \(\mathbf{C}_n\) 的长度。根据命题 2 a)，序列 \(d_n\)（其元素属于 \(\mathbf{N} \cup \{+\infty\}\)）是递增的；若它是平稳的，则存在整数 \(n_0\)，使得 \(\pi_{n,m}\) 是从 \(\mathbf{C}_m\) 到 \(\mathbf{C}_n\) 的同构，当 \(n_0 \leqslant n \leqslant m\) 时，于是 C 同构于 \(\mathbf{C}_{n_0}\)，是一个 \(p\)-环。

因此只需考虑每个 \(d_{n}\) 都有限且序列 \((d_{n})\) 趋于 \(+\infty\) 的情形。赋予环 C 平凡滤子（III，§ 2，第 1 号，例 5）。对于 \(n\in\mathbf{N}\)，令 \(\mathbf{I}_{n}\) 为 \(\pi_{n}\) 的核；置 \(\mathbf{I}_{n}=\mathbf{C}\)（当 \(n<0\) 时）。令 E 表示带有滤子 \((\mathbf{I}_{n})_{n\in\mathbf{Z}}\) 的 C-模 C。它是分离且完备的，因为拓扑 \(\mathcal{G}\) 由滤子 \((\mathbf{I}_{n})_{n\in\mathbf{Z}}\) 定义，正是各 \(\mathbf{C}_{n}\) 上离散拓扑的射影极限拓扑。

令 \(k\) 为满足 \(\geqslant 1\) 的整数。有 \(p^{k}\mathrm{C} \subset \varprojlim (p^{k}\mathrm{C}_{n})\)（E，III，p. 55，公式（9））。反之，若 \(x = (x_{n})_{n \in \mathbb{N}} \in \varprojlim (p^{k}\mathrm{C}_{n})\)，并置 \(\mathrm{X}_{n} = \{y \in \mathrm{C}|\pi_{n}(p^{k}y) = x_{n}\}\)，则序列 \((\mathrm{X}_{n})_{n \in \mathbb{N}}\) 是 E 的非空闭仿射子集的递减序列。由于 \(\mathrm{E / I}_{n}\) 是阿廷 C-模，\(\mathrm{X}_{n}\) 的交非空（III，§ 2，第 7 号，命题 7）；对于每个 \(z \in \bigcap_{n \in \mathbb{N}} \mathrm{X}_{n}\)，都有 \(p^{k}z = x\)。因此我们证明了 \(p^{k}\mathrm{C} = \varprojlim p^{k}\mathrm{C}_{n}\)。

对于每个整数 \(k \geqslant 1\)。特别地，C 的理想 \(p^{k}\text{C}\) 关于拓扑 \(\mathcal{G}\) 是闭的。在 C 上，\(p\)-进拓扑比拓扑 \(\mathcal{G}\) 更细，因为有 \(p^{d_{n}}\text{C} \subset \text{I}_{n}\)。于是根据 TG，III，p. 26，命题 10 的推论 1，C 关于 \(p\text{C}\)-进拓扑是分离且完备的。此外，有 \(p\text{C} = \varprojlim p\text{C}_{n} = \pi_{0}^{-1}(p\text{C}_{0})\)，所以从 C/pC 到 \(\text{C}_{0}/p\text{C}_{0}\) 的满射同态（由 \(\pi_{0}\) 得到）是同构。这说明 C 的理想 \(p\text{C}\) 是极大理想，因而 C 是一个 \(p\)-环。命题 3 的最后一个断言由命题 2 b) 得出。

